\begin{homeworkProblem}
  Seja $G$ um grupo que tem precisamente dois subgrupos. Mostre que $G$ é finito de ordem prima.
  \begin{solution}
    Como $G$ possui precisamente dois subgrupos, os únicos subgrupos de ele são $\{1\}$ e $G$.\\
    Vamos a raciocinar por contradição. Então suponha que $G$ no é finito, logo $\langle x\rangle \neq \langle x^{2}\rangle$ isso porque $\langle x\rangle = G$ e $o(x)=\infty$, mas isso é uma contradição porque $\{1\}$ e $G$ são os únicos subgrupos de $G$. Então podemos afirmar que $G$ é finito.
    Por outro lado agora suponha que $|G|=m$ com $m$ não prima, então sem perdida de generalidade suponha $m=p_1p_2$ com $p_1$ e $p_2$ primas, logo tome $x\in G$ tal que $x\neq 1$, logo como $\langle x\rangle = G$ voce pode tomar $y=x^{p_1}$. Então nos sabemos que $\langle y\rangle \neq \{1\}$ e $\langle y\rangle\neq G$, isso e uma contradição porque $\{1\}$ e $G$ são os únicos subgrupos de $G$. Logo voce pode concluir que si $G$ é um grupo que tem precisamente dois subgrupos, então $G$ é finito de ordem prima. 
  \end{solution}
\end{homeworkProblem}
